% ======================================================================
% capitulos/capitulo5.tex
% ----------------------------------------------------------------------
% ENTREGABLE 5 — SOLUCIÓN TIC INTEGRADA al Ecosistema de Sistemas de Información
% Empresa caso de estudio: WISP Ingeniería en Redes y Telecomunicaciones
% Stack: Django · Angular · PostgreSQL · App móvil · WireGuard · API RouterOS
% Contenido trasladado desde Entregable5_SolucionTIC_GTI_WISP.docx
% ======================================================================
\chapter{Solución TIC Integrada al Ecosistema de Sistemas de Información}
\label{chap:soluciontic}

% ----------------------------------------------------------------------
% Descripción del producto (introducción del capítulo)
% ----------------------------------------------------------------------
El presente documento constituye el diseño conceptual y técnico de la solución
tecnológica propuesta para WISP: el \textbf{Sistema de Gestión Integral para ISP de
Fibra Óptica Multi-Zona}. Considera su integración con los sistemas y la
infraestructura existentes, la arquitectura organizacional, los flujos de datos y los
requisitos de seguridad.

Se describe la arquitectura de la solución y sus componentes, el ecosistema de
sistemas de información y su interoperabilidad, los mecanismos de seguridad y
gobernanza (alineados a COBIT 2019 y a la normativa peruana de protección de datos),
y los lineamientos de escalabilidad y sostenibilidad, asegurando la coherencia con el
ecosistema institucional y con los objetivos estratégicos definidos en los
entregables previos.

\begin{upeubox}{Fundamento metodológico: descripción de arquitectura}
La arquitectura se documenta siguiendo el enfoque de vistas del modelo C4 (contexto,
contenedores, componentes) y el estándar ISO/IEC/IEEE~42010 de descripción de
arquitectura, complementados con un modelo de capas (\textit{layered architecture})
para representar la separación de responsabilidades.
\end{upeubox}

% ======================================================================
\section{Arquitectura de la Solución}
\label{sec:arquitectura}

La solución adopta una arquitectura cliente-servidor multicapa, desplegada en la nube
y conectada de forma segura con la infraestructura física distribuida en las 13 zonas
de WISP. El backend centralizado expone una API REST que consumen tanto el frontend
web como la aplicación móvil, y se comunica con los 13 MikroTik a través de una red
privada virtual (WireGuard) usando la API nativa de RouterOS~7.

\subsection{Diagrama de arquitectura (vista de despliegue)}

\begin{center}
\includegraphics[width=\textwidth]{arq_despliegue.png}
\captionof{figure}{Arquitectura de despliegue: componentes en la nube (VPS) conectados
vía WireGuard con los 13 MikroTik y OLT de cada zona.}
\label{fig:arq-despliegue}
\end{center}

\subsection{Arquitectura por capas}

La solución se organiza en cinco capas con responsabilidades bien separadas, lo que
favorece el mantenimiento, la seguridad y la escalabilidad:

\begin{center}
\includegraphics[width=0.9\textwidth]{arq_capas.png}
\captionof{figure}{Arquitectura en capas (\textit{layered architecture}) de la
solución.}
\label{fig:arq-capas}
\end{center}

\subsection{Componentes principales}

\begin{center}
\captionof{table}{Componentes principales de la solución.}
\label{tab:componentes}
\begin{tabularx}{\textwidth}{L{3cm} L{3.2cm} X}
\toprule
\textbf{Componente} & \textbf{Tecnología} & \textbf{Función} \\
\midrule
Frontend web & Angular &
Panel de administración: gestión de clientes, planes, zonas, órdenes y cobertura. \\
Aplicación móvil & App móvil (roles) &
Registro de prospectos (vendedor) y evidencias de instalación (técnico) desde el celular. \\
Backend / API & Django REST Framework &
Lógica de negocio, autenticación por roles, exposición de API REST y orquestación de
integraciones. \\
Base de datos & PostgreSQL &
Fuente única de verdad: clientes, planes, zonas, órdenes, instalaciones, usuarios. \\
VPN & WireGuard &
Túnel cifrado entre la nube y los 13 MikroTik para gestión remota segura. \\
Integración de red & API RouterOS~7 &
Alta, corte y reactivación automática de usuarios PPPoE en los MikroTik. \\
Notificaciones & Node.js (WhatsApp) &
Envío automático de mensajes a clientes (avisos, cobros, estados). \\
\bottomrule
\end{tabularx}
\end{center}

\subsection{Integraciones}

La solución integra tres tipos de sistemas, lo que constituye su principal
diferenciador técnico:

\begin{itemize}[leftmargin=1.4em]
  \item \textbf{Integración con MikroTik (RouterOS API):} el backend crea, corta y
  reactiva usuarios PPPoE y gestiona la lista de corte de morosos directamente en cada
  router, eliminando la operación manual.
  \item \textbf{Integración con WhatsApp (Node.js):} envío automatizado de
  notificaciones a los clientes, reutilizando el canal ya conocido por la empresa.
  \item \textbf{Integración interna (web $+$ móvil $+$ BD):} todos los clientes (web y
  móvil) consumen la misma API y la misma base de datos, garantizando consistencia.
\end{itemize}

% ======================================================================
\section{Ecosistema de Sistemas de Información}
\label{sec:ecosistema}

\subsection{Mapa de sistemas actuales}

El ecosistema actual de WISP está compuesto por sistemas aislados que no se comunican
entre sí, lo que genera la dispersión de datos identificada en el diagnóstico:

\begin{center}
\captionof{table}{Mapa de sistemas actuales y su destino en la solución.}
\label{tab:mapa-sistemas}
\begin{tabularx}{\textwidth}{L{3.6cm} X X}
\toprule
\textbf{Sistema actual} & \textbf{Rol} & \textbf{Destino en la solución} \\
\midrule
Sistema heredado (PHP) & Gestión de clientes & Reemplazado por el nuevo sistema web \\
13$\times$ MikroTik RouterOS~7 & Control PPPoE y ancho de banda & Integrado vía API \\
OLT VSOL ($\times$13) & Concentración de fibra (GPON) & Se mantiene (capa física) \\
Bot WhatsApp (pagos) & Recepción de comprobantes & Coexiste (fuera del alcance actual) \\
WhatsApp operativo & Coordinación informal & Reemplazado por app $+$ notificaciones \\
Papel / registros manuales & Captación en campo & Reemplazado por app móvil \\
\bottomrule
\end{tabularx}
\end{center}

\subsection{Integración propuesta}

La solución actúa como núcleo integrador del ecosistema: centraliza la información
dispersa y se convierte en el punto único de gestión. El sistema heredado y los flujos
de papel/WhatsApp se sustituyen; los MikroTik y OLT (infraestructura física) se
conservan pero pasan a ser gestionados desde el sistema; y el canal WhatsApp se
automatiza. De un ecosistema fragmentado se pasa a un ecosistema integrado alrededor de
una fuente única de verdad.

\begin{upeubox}{Principio de integración adoptado}
Se aplica el principio de «fuente única de verdad» (\textit{single source of truth}):
toda la información operativa reside en la base de datos PostgreSQL, y los distintos
clientes (web, móvil) y sistemas (MikroTik, WhatsApp) se sincronizan con ella. Esto
elimina la duplicidad y las inconsistencias que hoy existen entre el sistema heredado,
los comentarios de los MikroTik y los chats.
\end{upeubox}

\subsection{Interoperabilidad}

La interoperabilidad entre los componentes se garantiza mediante estándares abiertos y
protocolos ampliamente soportados:

\begin{center}
\captionof{table}{Interoperabilidad entre componentes de la solución.}
\label{tab:interoperabilidad}
\begin{tabularx}{\textwidth}{L{3.8cm} X L{3cm}}
\toprule
\textbf{Interacción} & \textbf{Protocolo / Estándar} & \textbf{Formato} \\
\midrule
Web/Móvil $\leftrightarrow$ Backend & HTTP/HTTPS (REST) & JSON \\
Backend $\leftrightarrow$ MikroTik & API REST RouterOS~7 sobre VPN & JSON / HTTPS \\
Nube $\leftrightarrow$ Zonas & WireGuard (UDP cifrado) & Túnel VPN \\
Backend $\leftrightarrow$ WhatsApp & API / librería Node.js & JSON \\
Backend $\leftrightarrow$ Base de datos & SQL / ORM (Django) & Relacional \\
\bottomrule
\end{tabularx}
\end{center}

% ======================================================================
\section{Seguridad y Gobernanza}
\label{sec:seguridad}

La seguridad y la gobernanza de la solución se abordan en dos planos complementarios:
los controles técnicos de seguridad de la información, y el marco de gobierno de TI que
asegura que la tecnología esté alineada con los objetivos de WISP y sea gestionada con
criterios reconocidos.

\begin{upeubox}{Fundamento metodológico: gobierno de TI (COBIT 2019)}
La gobernanza se alinea con COBIT~2019 (ISACA), marco de referencia para el gobierno y
la gestión de la tecnología de la información empresarial. Se referencian objetivos de
gobierno y gestión (\textit{governance/management objectives}) aplicables a la
solución. Esto responde al lineamiento de «Gobierno de TI» que el propio programa de
GTI establece para el área.
\end{upeubox}

\subsection{Gestión de accesos}

El control de accesos se basa en el principio de mínimo privilegio y en la autenticación
por roles (RBAC --- \textit{Role-Based Access Control}), de modo que cada usuario accede
únicamente a las funciones que su rol requiere:

\begin{center}
\captionof{table}{Control de accesos por rol (RBAC).}
\label{tab:accesos}
\begin{tabularx}{\textwidth}{L{2.6cm} X L{3.6cm}}
\toprule
\textbf{Rol} & \textbf{Accesos permitidos} & \textbf{Restricciones} \\
\midrule
Administrador &
Gestión completa: clientes, planes, zonas, órdenes, cortes, reportes. &
Acceso total (supervisado) \\
Vendedor &
Registro de prospectos y consulta de cobertura desde la app. &
No accede a datos de red ni cortes \\
Técnico &
Órdenes asignadas y registro de evidencias de instalación. &
No accede a datos comerciales/financieros \\
Cliente (futuro) &
Consulta de su estado y cobertura (autoservicio). &
Solo su propia información \\
\bottomrule
\end{tabularx}
\end{center}

\subsection{Protección de datos}

La protección de los datos ---especialmente los datos personales de los clientes--- se
garantiza mediante los siguientes controles técnicos:

\begin{itemize}[leftmargin=1.4em]
  \item \textbf{Cifrado en tránsito:} todo el tráfico web usa HTTPS/TLS, y la
  comunicación con los MikroTik viaja por el túnel cifrado WireGuard.
  \item \textbf{Autenticación robusta:} credenciales protegidas (\textit{hashing}),
  tokens de sesión y control por roles.
  \item \textbf{Respaldos periódicos:} copias de seguridad automáticas de la base de
  datos, con plan de recuperación.
  \item \textbf{Registro de auditoría:} trazabilidad de las acciones críticas (altas,
  cortes, cambios) para responsabilidad y control.
  \item \textbf{Segmentación de red:} acceso a los MikroTik restringido a través de la
  VPN, sin exposición directa a internet.
\end{itemize}

\subsection{Cumplimiento normativo}

La solución se diseña considerando el marco legal peruano y las buenas prácticas del
sector:

\begin{center}
\captionof{table}{Marco normativo aplicable a la solución.}
\label{tab:normativo}
\begin{tabularx}{\textwidth}{L{3cm} L{3.4cm} X}
\toprule
\textbf{Norma / Marco} & \textbf{Ámbito} & \textbf{Aplicación en la solución} \\
\midrule
Ley N° 29733 & Protección de Datos Personales (Perú) &
Tratamiento seguro y consentido de los datos de clientes. \\
OSIPTEL / MTC & Regulación de telecomunicaciones &
Registro ordenado de clientes y servicios prestados. \\
ISO/IEC 27001 (ref.) & Seguridad de la información &
Buenas prácticas de controles de seguridad. \\
COBIT 2019 & Gobierno y gestión de TI &
Alineamiento estratégico y gestión de riesgos (ver abajo). \\
\bottomrule
\end{tabularx}
\end{center}

\subsection{Gobernanza de TI con COBIT 2019}

Se referencian los objetivos de gobierno y gestión de COBIT~2019 que la solución ayuda
a cumplir, evidenciando que el proyecto no solo construye tecnología, sino que la
gobierna de forma alineada al negocio:

\begin{center}
\captionof{table}{Objetivos de COBIT~2019 cubiertos por la solución.}
\label{tab:cobit}
\begin{tabularx}{\textwidth}{L{3.4cm} L{3.6cm} X}
\toprule
\textbf{Objetivo COBIT 2019} & \textbf{Dominio} & \textbf{Contribución de la solución} \\
\midrule
EDM01 -- Gobierno de TI & Evaluar, Orientar y Supervisar &
La solución alinea la TI con los objetivos estratégicos de WISP definidos en el
diagnóstico. \\
EDM03 -- Optimización del riesgo & Evaluar, Orientar y Supervisar &
Reduce riesgos operativos (dependencia de una persona, pérdida de datos). \\
APO12 -- Gestión del riesgo & Alinear, Planificar y Organizar &
Se apoya en el plan de riesgos del proyecto (análisis P$\times$I y respuestas). \\
APO13 -- Gestión de la seguridad & Alinear, Planificar y Organizar &
Define políticas de seguridad, cifrado y control de accesos. \\
BAI03 -- Gestión de soluciones & Construir, Adquirir e Implementar &
Desarrollo estructurado del sistema (SRS, arquitectura, pruebas). \\
DSS05 -- Servicios de seguridad & Entregar, Dar Servicio y Soporte &
Gestión de accesos por roles, protección de \textit{endpoints} y VPN. \\
\bottomrule
\end{tabularx}
\end{center}

\begin{upeubox}{Lectura de gobernanza}
Al mapear la solución contra COBIT~2019, WISP no solo obtiene un sistema, sino un
modelo de gobierno de su TI: la tecnología queda alineada con la estrategia (EDM01),
los riesgos se gestionan (EDM03, APO12), la seguridad se formaliza (APO13, DSS05) y el
desarrollo sigue buenas prácticas (BAI03). Esto eleva la madurez digital de la empresa
desde el nivel «inicial» diagnosticado hacia un nivel «gestionado».
\end{upeubox}

% ======================================================================
\section{Escalabilidad y Sostenibilidad}
\label{sec:escalabilidad}

\subsection{Capacidad de crecimiento}

La arquitectura está diseñada para crecer sin rediseños estructurales, acompañando la
expansión de WISP a nuevos distritos (objetivo estratégico OE-01):

\begin{itemize}[leftmargin=1.4em]
  \item \textbf{Escalabilidad horizontal de zonas:} agregar una nueva zona solo
  requiere sumar su MikroTik a la VPN y registrarlo; el sistema ya está preparado para
  $N$ zonas, no está limitado a 13.
  \item \textbf{Escalabilidad de usuarios:} el backend Django y PostgreSQL soportan el
  crecimiento de clientes; ante mayor carga, el VPS puede escalar verticalmente (más
  recursos) u horizontalmente (más instancias).
  \item \textbf{Homogeneidad de la infraestructura:} al ser todos los MikroTik del
  mismo modelo y versión (RB4011, RouterOS~7), la integración escala sin casos
  particulares.
  \item \textbf{Arquitectura desacoplada:} la separación en capas y el uso de API
  permiten evolucionar cada componente de forma independiente.
\end{itemize}

\subsection{Costos futuros}

La sostenibilidad económica se sustenta en costos operativos bajos y predecibles, que
crecen mucho más lento que los ingresos por nuevos clientes:

\begin{center}
\captionof{table}{Proyección de costos operativos según escenario de crecimiento.}
\label{tab:costos-futuros}
\begin{tabularx}{\textwidth}{L{4.2cm} L{3.4cm} X}
\toprule
\textbf{Escenario} & \textbf{Costo operativo anual} & \textbf{Observación} \\
\midrule
Actual ($\sim$455 clientes, 13 zonas) & $\sim$S/~3\,780 &
VPS $+$ dominio $+$ mantenimiento $+$ backup \\
Crecimiento ($+$50\% clientes) & $\sim$S/~4\,500--5\,500 &
Solo escala el VPS; el resto se mantiene \\
Nuevas zonas ($+$5 zonas) & Costo marginal bajo &
Sin nuevo desarrollo; solo configuración \\
\bottomrule
\end{tabularx}
\end{center}

\begin{upeubox}{Sostenibilidad económica}
El costo operativo por cliente disminuye a medida que WISP crece (economía de escala):
la misma plataforma soporta más clientes y zonas con un incremento de costo marginal.
Esto hace que la solución sea sostenible en el tiempo y refuerza el retorno demostrado
en el Business Case.
\end{upeubox}

\subsection{Evolución tecnológica}

La solución contempla una hoja de ruta de evolución que aprovecha su arquitectura
desacoplada para incorporar mejoras futuras sin rehacer lo existente:

\begin{center}
\captionof{table}{Hoja de ruta de evolución tecnológica.}
\label{tab:evolucion}
\begin{tabularx}{\textwidth}{L{2.8cm} X L{2.4cm}}
\toprule
\textbf{Fase} & \textbf{Evolución propuesta} & \textbf{Horizonte} \\
\midrule
Corto plazo &
Automatización del módulo de pagos (Yape/Plin) e integración con el bot de cobros. &
6--12 meses \\
Mediano plazo &
Portal de autoservicio para clientes y provisión automática de ONUs en la OLT. &
1--2 años \\
Largo plazo &
Analítica de datos e indicadores de negocio (BI); alertas predictivas de red. &
2$+$ años \\
\bottomrule
\end{tabularx}
\end{center}

% ======================================================================
\section{Conclusión — Cierre de la Línea GTI}
\label{sec:cierre-gti}

La Solución TIC Integrada propuesta constituye una respuesta coherente, técnicamente
sólida y estratégicamente alineada al problema de WISP. Su arquitectura multicapa en la
nube, integrada de forma segura con los 13 MikroTik mediante WireGuard y la API de
RouterOS, convierte un ecosistema hoy fragmentado en una plataforma única, trazable y
escalable.

Con este documento se cierra la línea de Gestión de Tecnologías de Información (GTI),
que ha recorrido de forma completa el ciclo: el diagnóstico organizacional identificó el
problema; el Business Case demostró su viabilidad económica; el Plan de Gestión
estableció cómo ejecutarlo; el modelado de procesos AS-IS/TO-BE cuantificó la mejora
operativa; y esta Solución TIC define el cómo técnico. En conjunto, GTI entrega el marco
estratégico y de gobierno sobre el cual las líneas de Software e Infraestructura
desarrollarán la implementación.

\begin{upeubox}{Coherencia integral del proyecto}
Los cinco componentes de GTI conforman una narrativa única y trazable: problema
\textrightarrow{} viabilidad \textrightarrow{} plan \textrightarrow{} procesos
\textrightarrow{} solución. Cada decisión técnica de esta Solución TIC responde a una
necesidad identificada en el diagnóstico y validada en el Business Case, y se gobierna
con criterios reconocidos (PMBOK, BPMN, COBIT~2019). El proyecto queda listo para su
fase de construcción en las líneas de Software e Infraestructura.
\end{upeubox}
